\chapter{Estado del arte}
\label{ch:estado-arte}

La optimización de carteras es uno de los problemas más estudiados de las
finanzas cuantitativas. Desde su formulación original a mediados del siglo~XX,
la disciplina ha evolucionado en torno a una tensión persistente: la elegancia
del modelo teórico frente a la fragilidad de su aplicación práctica, dominada
por el error con que se estiman sus parámetros. Este capítulo revisa esa
evolución ---de la teoría moderna de carteras a la irrupción del aprendizaje
automático--- con el doble objetivo de situar las aportaciones más relevantes
de la literatura y de delimitar el espacio que ocupa este trabajo.

%--------------------------------------------------------------------
\section{De la teoría moderna de carteras al CAPM y la APT}
%--------------------------------------------------------------------

El punto de partida es el trabajo seminal de \citet{Markowitz1952}, que
formalizó la selección de carteras como un compromiso entre la rentabilidad
esperada y el riesgo, medido este último por la varianza. Su principal
aportación conceptual fue demostrar que el riesgo de una cartera no es la suma
de los riesgos individuales, sino que depende crucialmente de las covarianzas
entre activos: la diversificación permite reducir el riesgo sin sacrificar
rentabilidad esperada. Este marco media-varianza constituye todavía hoy el
lenguaje básico de la gestión cuantitativa.

Poco después, \citet{Tobin1958} introdujo el \textit{teorema de separación}: al
incorporar un activo libre de riesgo, la decisión de inversión se descompone en
dos etapas independientes ---la elección de una única cartera óptima de activos
con riesgo (la cartera tangente) y su combinación con el activo seguro---. Esta
idea pavimentó el camino hacia los modelos de equilibrio. El \textit{Capital
Asset Pricing Model} (CAPM), desarrollado de forma independiente por
\citet{Sharpe1964} y \citet{Lintner1965}, postuló que, en equilibrio, la
rentabilidad esperada de un activo es una función lineal de su riesgo
sistemático (la beta respecto al mercado). Como alternativa más general,
\citet{Ross1976} propuso la \textit{Arbitrage Pricing Theory} (APT), que
explica las rentabilidades mediante varios factores de riesgo sistemáticos en
lugar de uno solo, sin apoyarse en los supuestos de equilibrio del CAPM.

En el plano matemático, \citet{Merton1972} obtuvo la solución analítica
cerrada de la frontera eficiente en ausencia de restricciones de desigualdad,
mostrando que el conjunto de carteras de mínima varianza describe una hipérbola
en el plano riesgo-rentabilidad. Este resultado, que se retoma con detalle en
el Capítulo~\ref{ch:formulacion}, es la base sobre la que se construye toda la
maquinaria de optimización del presente trabajo.

%--------------------------------------------------------------------
\section{El problema del error de estimación}
%--------------------------------------------------------------------

La elegancia del modelo de Markowitz contrasta con su comportamiento práctico.
El modelo requiere como entradas el vector de rentabilidades esperadas y la
matriz de covarianzas, parámetros que en la práctica no se conocen y deben
estimarse a partir de datos históricos. \citet{Merton1980} fue uno de los
primeros en señalar que estimar la rentabilidad esperada es mucho más difícil
que estimar las varianzas y covarianzas: mientras que la precisión de estas
últimas mejora al aumentar la frecuencia de muestreo, la de las primeras solo
mejora ampliando el horizonte temporal, lo que exige series de datos
inviablemente largas.

Las consecuencias de este hecho son severas. \citet{BestGrauer1991} demostraron
analíticamente que las carteras eficientes son extremadamente sensibles a
pequeñas variaciones en las rentabilidades esperadas estimadas: cambios mínimos
en las medias producen reasignaciones drásticas de los pesos. En la misma línea,
\citet{Michaud1989} acuñó la célebre expresión de que la optimización
media-varianza actúa como un \textit{«maximizador del error de estimación»},
ya que tiende a asignar pesos desproporcionados precisamente a aquellos activos
cuya rentabilidad ha sido sobrestimada por azar; de ahí que las carteras
``optimizadas'' a menudo rindan peor, fuera de muestra, que alternativas
ingenuas.

La evidencia empírica más influyente en este sentido es la de
\citet{DeMiguel2009}, quienes compararon catorce modelos de optimización frente
a la sencilla regla de diversificación ingenua $1/N$ (repartir el capital por
igual entre todos los activos). Su conclusión fue contundente: ninguna de las
estrategias optimizadas superaba consistentemente a la $1/N$ fuera de muestra, y
estimaron que se necesitarían ventanas de estimación de miles de meses para que
la ventaja teórica de la optimización compensara el error de estimación. Este
resultado es el referente directo frente al que se contrastan los hallazgos del
presente trabajo.

%--------------------------------------------------------------------
\section{Mejoras en la estimación de la matriz de covarianza}
%--------------------------------------------------------------------

Aunque la covarianza es más fácil de estimar que la rentabilidad esperada,
también sufre cuando el número de activos es comparable al de observaciones: la
matriz de covarianza muestral se vuelve mal condicionada y produce carteras
inestables. Una de las soluciones más exitosas es el \textit{shrinkage} de
\citet{Ledoit2004}, que estima la covarianza como una combinación convexa, de
intensidad óptima, entre la matriz muestral y un objetivo bien condicionado (la
matriz identidad escalada). El estimador resultante está mejor condicionado,
reduce el error cuadrático medio y da lugar a carteras más estables y
diversificadas.

Una vía complementaria es la imposición de una estructura de factores.
\citet{FamaFrench1993} mostraron que un reducido número de factores comunes
---el mercado, el tamaño y el valor--- explica buena parte de la variación
conjunta de las rentabilidades de las acciones. Modelar la covarianza a través
de pocos factores reduce drásticamente el número de parámetros a estimar y
filtra el ruido idiosincrásico; esta idea enlaza con la APT de
\citet{Ross1976} y con la descomposición en componentes principales que se
emplea en este trabajo como tercer estimador de la covarianza
(Capítulo~\ref{ch:limitaciones}).

%--------------------------------------------------------------------
\section{Enfoques bayesianos e incorporación de visiones}
%--------------------------------------------------------------------

Una familia distinta de soluciones aborda directamente el error en las
rentabilidades esperadas mediante la inferencia bayesiana. \citet{Jorion1986}
aplicó el estimador de Bayes-Stein, que contrae las rentabilidades esperadas
individuales hacia un valor común (típicamente el de la cartera de mínima
varianza), reduciendo la varianza del estimador y mejorando el rendimiento
fuera de muestra. El enfoque más conocido en la práctica profesional es, no
obstante, el modelo de \citet{BlackLitterman1992}, que parte de las
rentabilidades implícitas en el equilibrio de mercado y las combina, en un marco
bayesiano, con las ``visiones'' subjetivas del inversor. El resultado son
carteras más estables e intuitivas que las que produce la optimización
media-varianza directa, lo que mitiga el problema de las asignaciones extremas.
Estos enfoques son conceptualmente afines a la sustitución del estimador ingenuo
de la rentabilidad esperada por una predicción mejor informada, que es
precisamente lo que en este trabajo se aborda mediante el aprendizaje
automático.

%--------------------------------------------------------------------
\section{Predicción de rentabilidades y aprendizaje automático}
%--------------------------------------------------------------------

La pregunta de fondo ---¿es posible predecir las rentabilidades?--- tiene una
respuesta matizada. Por un lado, la literatura ha documentado anomalías
empíricas explotables, como el efecto \textit{momentum} de \citet{Jegadeesh1993},
según el cual los activos que han obtenido buenas rentabilidades en el pasado
reciente tienden a seguir haciéndolo a corto plazo. Por otro,
\citet{Cont2001} sistematizó los \textit{hechos estilizados} de las series
financieras ---colas pesadas, ausencia de autocorrelación lineal en los
rendimientos y agrupamiento de la volatilidad---, que limitan severamente la
predictibilidad de la rentabilidad y advierten contra los supuestos de
normalidad.

En la última década, el aprendizaje automático ha irrumpido con fuerza en este
terreno. El estudio de referencia es el de \citet{Gu2020}, que evaluó de forma
sistemática un amplio catálogo de métodos ---desde regresiones penalizadas hasta
árboles y redes neuronales--- para la predicción de rentabilidades en un gran
panel de acciones. Sus conclusiones son especialmente relevantes para este
trabajo: los métodos no lineales pueden capturar interacciones que los lineales
ignoran, pero las ganancias predictivas son modestas (con $R^2$ fuera de muestra
muy pequeños), y la regularización resulta decisiva para evitar el sobreajuste
en un entorno de baja señal-ruido. Entre los métodos concretos, este trabajo
emplea los \textit{Random Forests} de \citet{Breiman2001} ---un modelo no
lineal robusto basado en el promediado de árboles--- y dos regresiones
regularizadas: el \textit{Ridge} y el \textit{Lasso}. Este último, emparentado
con la \textit{elastic net} de \citet{Zou2005}, realiza además selección
automática de variables, lo que resulta especialmente útil en entornos de baja
señal-ruido. Estas predicciones de la rentabilidad esperada son las que
alimentan el optimizador.

Conviene subrayar que la incorporación de estas predicciones al modelo de
Markowitz preserva su naturaleza de problema de optimización convexa, que puede
resolverse de forma exacta con las técnicas descritas por
\citet{Boyd2004}, sin recurrir a heurísticas. Esta combinación de predicción
estadística y optimización exacta es el eje metodológico del presente trabajo.

%--------------------------------------------------------------------
\section{Evaluación del rendimiento fuera de muestra}
%--------------------------------------------------------------------

La medida estándar para comparar estrategias de inversión es el ratio de
\citet{Sharpe1966}, que cuantifica la rentabilidad obtenida por unidad de riesgo
asumido. Sin embargo, comparar ratios de Sharpe a partir de una muestra finita
exige cautela estadística: las diferencias observadas pueden deberse al azar.
\citet{Memmel2003} corrigió el contraste de Jobson-Korkie para obtener un test
válido de la igualdad de dos ratios de Sharpe, herramienta que este trabajo
emplea para evaluar la significancia de sus resultados. La propia metodología de
comparación fuera de muestra frente a baselines ingenuos, popularizada por
\citet{DeMiguel2009}, constituye el estándar metodológico que aquí se adopta.

%! Creo que la palabra significancia no aparece contemplada por la RAE. Utiliza en su lugar significación.

%--------------------------------------------------------------------
\section{Posicionamiento de este trabajo}
%--------------------------------------------------------------------

A la luz de la literatura revisada, este trabajo se sitúa en la intersección de
tres líneas: la formulación matemática rigurosa del problema de Markowitz como
programa cuadrático convexo resuelto de forma exacta \cite{Markowitz1952,
Boyd2004}; la mitigación del error de estimación tanto en la covarianza
---mediante \textit{shrinkage} y un modelo de factores por componentes
principales \cite{Ledoit2004, Ross1976}--- como en la rentabilidad esperada
---mediante aprendizaje automático \cite{Gu2020}---; y la evaluación honesta del rendimiento fuera de
muestra con contraste de significancia estadística \cite{DeMiguel2009,
Memmel2003}.

Frente a la abundante literatura de carácter aplicado, que con frecuencia
recurre a métodos heurísticos de optimización y a un enfoque orientado al
negocio, la aportación diferencial de este trabajo es doble. En primer lugar,
el énfasis en el \textbf{rigor matemático}: la resolución exacta del problema
convexo, el análisis de sus condiciones de optimalidad (KKT) y de su dualidad, y
la conexión explícita entre la descomposición espectral de la covarianza y la
teoría de factores. En segundo lugar, la \textbf{honestidad empírica}: lejos de
presentar una estrategia ganadora, el trabajo somete sus resultados a un
contraste de significancia estadística y discute con transparencia las
condiciones bajo las cuales la optimización ayuda ---y aquellas, más frecuentes,
en las que el error de estimación impide distinguirla de una diversificación
ingenua.
